\section{Conclusion}
\label{sec:conclusion}

The reviewed evidence answers the research question as follows: learned
instruction--data separation can substantially reduce prompt-injection success
under the attacker models represented in the original StruQ and SecAlign
evaluations. SecAlign provides stronger resistance than StruQ against the
optimisation-based attacks included in its study
\parencites{chen2025struq}{chen2025secalign}. However, broader and
defence-adaptive white-box evaluations show that neither approach provides an
independently enforced security boundary
\parencites{jia2025critical}
             {yang2025checkpointgcg}
             {pandya2025astra}.

Their effectiveness therefore depends on the evaluated model, task, attack
construction method, attacker access, optimisation budget, and success
criterion. Strong results against handcrafted attacks or a fixed target output
should not be generalised to stronger adaptive attackers or different
application settings. The later evaluations indicate that some early results
were partly enabled by narrow test scopes and attack suites that were not fully
adapted to the defence
\parencite{jia2025critical}.

Learned instruction--data separation should consequently be treated as one
layer of risk reduction rather than complete prevention. Applications must
still determine which inputs are trusted and must restrict what the model can
access or cause. Least privilege, external authorisation, containment,
monitoring, and independent approval are required when model outputs can expose
protected information or initiate consequential actions
\parencites{owasp2025llm}
             {wu2024systemlevel}
             {debenedetti2024agentdojo}
             {debenedetti2025defeating}.

The available evidence remains limited by differences in models, datasets,
tasks, attacker access, and evaluation methods. Further research should compare
defences under shared benchmarks, defence-adaptive attacks, realistic tool use,
and security outcomes such as disclosure or unauthorised actions rather than
target-string compliance alone. Until such evaluations provide stronger
evidence, prompting and learned instruction priority should not be treated as
security boundaries.